\documentclass[12pt,a4paper]{article}

\usepackage[margin=1in]{geometry}
\usepackage{times}
\usepackage{setspace}
\usepackage{graphicx}
\usepackage{array}
\usepackage{booktabs}
\usepackage{float}
\usepackage{fancyhdr}
\usepackage{titlesec}
\usepackage{hyperref}

\onehalfspacing

\setlength{\parindent}{0pt}
\setlength{\parskip}{6pt}

\pagestyle{fancy}
\fancyhf{}
\lhead{Campus Maintenance Intelligence}
\rhead{Project Proposal}
\cfoot{\thepage}

\begin{document}

% =========================
% FIRST PAGE
% =========================

\begin{titlepage}

\begin{center}

\vspace*{1.2cm}

{\Huge \textbf{Campus Maintenance Intelligence}}

\vspace{0.4cm}

{\Large An Intelligent Campus Maintenance Management System}

\vspace{1cm}

{\large Project Proposal for UCS503P}

\vspace{1cm}

\textbf{Submitted to:}

\vspace{0.2cm}

Dr. Raghav B Venkataramaiyer

\vspace{0.2cm}

Department of Computer Science and Engineering

\vspace{0.2cm}

Thapar Institute of Engineering and Technology

\vspace{1cm}

\textbf{Submitted By}

\vspace{0.4cm}

\begin{tabular}{ll}
Name & Roll Number \\
Gaurish Garg & 1024170019 \\
Sumit Singh & 1024170018 \\
Arav Garg & 1024170418
\end{tabular}

\vfill

\textbf{Academic Year 2026--27}

\end{center}

\end{titlepage}

% =========================
% CONTENTS
% =========================

\tableofcontents
\newpage

% =========================
% 1
% =========================

\section{Introduction}

Campus Maintenance Intelligence is a smart maintenance management system designed for educational campuses. It manages complaints related to electrical systems, plumbing, AC, housekeeping, masonry, painting, fire systems and other facilities.

Traditional maintenance systems generally process complaints using first-come-first-served methods. This may result in critical problems being handled after less important complaints.

The proposed system uses complaint information, location, severity, affected users and historical maintenance data to assign suitable priorities.

The system follows five major stages:

\begin{center}
\textbf{Structure $\rightarrow$ Cluster $\rightarrow$ Prioritize $\rightarrow$ Predict $\rightarrow$ Optimize}
\end{center}

This allows the system to move from simple complaint registration towards intelligent and proactive maintenance.


% =========================
% 2
% =========================

\section{Problem Statement}

Large campuses generate many maintenance complaints every day. Handling these complaints manually creates several problems:

\begin{itemize}
    \item Critical complaints may not receive immediate attention.
    \item Multiple users may report the same problem.
    \item Recurring failures are difficult to identify.
    \item Historical maintenance data is not fully utilized.
    \item Technician and maintenance resources may be allocated inefficiently.
\end{itemize}

Therefore, a system is required that can automatically prioritize complaints, identify duplicates, predict possible failures and assist in maintenance planning.


% =========================
% 3
% =========================

\section{Proposed Solution}

The proposed system provides a centralized platform for campus maintenance.

When a complaint is submitted, the system:

\begin{enumerate}
    \item Extracts important information from the complaint.
    \item Identifies similar or duplicate complaints.
    \item Assigns a priority based on severity, safety, location and affected users.
    \item Stores the complaint as an incident or work order.
    \item Uses historical data to identify possible future failures.
    \item Helps optimize maintenance resources.
\end{enumerate}

For example, an electrical short circuit in a hostel can receive a higher priority than a routine fan repair even if the fan complaint was submitted earlier.


% =========================
% 4
% =========================

\section{Project Objectives}

The main objectives are:

\begin{itemize}
    \item Develop a centralized campus maintenance platform.
    \item Automatically prioritize maintenance complaints.
    \item Detect duplicate and related complaints.
    \item Identify recurring maintenance problems.
    \item Predict possible future failures.
    \item Improve technician and resource allocation.
    \item Provide a dashboard for maintenance staff.
    \item Compare intelligent prioritization with first-come-first-served processing.
\end{itemize}


% =========================
% 5
% =========================

\section{System Overview}

The system consists of the following major modules:

\begin{table}[H]
\centering
\renewcommand{\arraystretch}{1.3}

\begin{tabular}{p{0.30\textwidth} p{0.55\textwidth}}
\toprule
\textbf{Module} & \textbf{Purpose} \\
\midrule

Complaint Management &
Register and track maintenance complaints. \\

Priority Engine &
Assign priority according to severity and campus context. \\

Duplicate Detection &
Identify complaints referring to the same problem. \\

Incident Management &
Combine related complaints into a single incident. \\

Predictive Maintenance &
Identify assets that may fail in the future. \\

Optimization &
Help allocate technicians and maintenance resources. \\

Dashboard &
Display complaints, priorities and maintenance statistics. \\

\bottomrule
\end{tabular}

\caption{Major System Modules}

\end{table}


% =========================
% 6
% =========================

\section{Technology Stack}

\begin{table}[H]
\centering
\renewcommand{\arraystretch}{1.3}

\begin{tabular}{p{0.35\textwidth} p{0.45\textwidth}}
\toprule
\textbf{Component} & \textbf{Technology} \\
\midrule
Frontend & Next.js, React.js, Tailwind CSS \\
Backend & FastAPI, Python \\
Database & PostgreSQL \\
Machine Learning & XGBoost, Random Forest \\
NLP & Sentence Transformers \\
Similarity Search & FAISS \\
Data Processing & Pandas, NumPy \\
Optimization & OR-Tools \\
Version Control & Git, GitHub \\
Deployment & Docker \\
\bottomrule
\end{tabular}

\caption{Technology Stack}

\end{table}


% =========================
% 7
% =========================

\section{Dataset}

The project will use complaint and historical maintenance data.

Complaint data will contain information such as:

\begin{itemize}
    \item Complaint description
    \item Category and problem type
    \item Building and location
    \item Timestamp
    \item Priority
    \item Incident ID
\end{itemize}

Historical maintenance data will contain information such as building details, maintenance work orders, repair costs, maintenance priority and labor hours.

A synthetic benchmark dataset will initially be used for testing. It will represent campus locations such as hostels, blocks, lecture theatres, the CSED department and the library.

The dataset will contain maintenance categories including electrical, plumbing, AC, carpentry, masonry, painting, housekeeping and fire systems.


% =========================
% 8
% =========================

\section{Development Methodology}

The project will be developed using an Agile approach.

The major stages are:

\begin{enumerate}
    \item Requirement analysis
    \item Dataset preparation
    \item Database and system design
    \item Backend development
    \item Frontend development
    \item Machine learning development
    \item System integration
    \item Testing
    \item Deployment
\end{enumerate}


% =========================
% 9
% =========================

\section{Expected Outcomes}

The project is expected to provide:

\begin{itemize}
    \item Faster identification of critical complaints.
    \item Reduction in duplicate complaint handling.
    \item Better maintenance prioritization.
    \item Identification of recurring problems.
    \item Early detection of possible failures.
    \item Better utilization of maintenance resources.
    \item A centralized maintenance dashboard.
\end{itemize}


% =========================
% 10
% =========================

\section{Future Scope}

The system can later be extended with:

\begin{itemize}
    \item IoT sensors for real-time infrastructure monitoring.
    \item Image-based detection of infrastructure damage.
    \item Mobile application for maintenance workers.
    \item Automated technician assignment.
    \item Advanced failure prediction.
    \item Integration with the actual university maintenance system.
\end{itemize}


% =========================
% 11
% =========================

\section{Conclusion}

Campus Maintenance Intelligence aims to transform traditional campus maintenance into a more intelligent and data-driven process.

The system combines complaint management, priority assignment, duplicate detection, predictive maintenance and resource optimization in a single platform.

By using historical and synthetic maintenance data, the system can be developed and evaluated before deployment in a real campus environment.

\begin{center}

\vspace{0.5cm}

\textbf{Campus Maintenance Intelligence}

\textit{Report. Prioritize. Predict. Optimize.}

\end{center}

\end{document}